% GENERATED FILE. Edit canonical lessons, then run npm run generate:books.
% canonical-source-hash: fnv1a64:0da553d860f5c543
% canonical-lessons: MR-C454-samet, MR-C454-gairsamaj, MR-C454-rusne, MR-C454-manavne, MR-C454-hirmusne

\chapter{Sulking and Making Up}
\label{ch:mr-sulking-and-making-up}
\hlchaptermodality{454}

\begin{chapteropening}
\textbf{By the end of this chapter:} \emph{I can talk about a sulk, a misunderstanding and making up.}

Complete the last lesson of chapter 454: I can talk about a sulk, a misunderstanding and making up.
\end{chapteropening}

\section[\texorpdfstring{\mr{समेट} (sameṭ)}{समेट (sameṭ)}]{\mr{समेट} (sameṭ) — a reconciliation}

\label{lesson:MR-C454-samet}



\begin{quote}
Before the new one: say the Marathi for to snap at, then the Marathi for to grumble.
\end{quote}

\subsection*{You'll want to know: \mr{समेट}}
\textbf{\mr{समेट}} — \emph{sameṭ} — \textquotedblleft{}a reconciliation\textquotedblright{}.

\begin{quote}
\mr{दोन} \mr{मित्रांमध्ये} \mr{शेवटी} \mr{समेट} \mr{झाला}.
\end{quote}

\begin{quote}
\emph{don mitrāmmadhye śevṭī sameṭ jhālā.}
\end{quote}

\begin{quote}
In the end the two friends made up.
\end{quote}

Masculine.

\mr{समेट} \mr{होणे}, to be reconciled; \mr{समेट} \mr{घडवणे}, to bring two sides together.

\begin{cousinweb}
\textbf{\mr{समेट}} (\emph{sameṭ}) is inherited: CDIAL lists the Marathi verb \emph{sameṭṇẽ}, \textquotedblleft{}to roll up\textquotedblright{}, under Sanskrit \emph{saṁveṣṭayati}, \textquotedblleft{}wraps up\textquotedblright{} (CDIAL 13026); the noun, a quarrel wound up and settled, is usually connected with that verb.
\end{cousinweb}

\begin{grammarlens}[title={What you've built}]
One more word for describing what you lived and felt.
\end{grammarlens}

\subsection*{Your turn}
\begin{itemize}
\raggedright
  \item \emph{Say it:} \emph{sameṭ}
  \item \emph{Say it:} \emph{sameṭ}, once more
  \item \emph{Say it:} \emph{don mitrāmmadhye śevṭī sameṭ jhālā.} — In the end the two friends made up.
  \item \emph{Recall:} say the Marathi for to snap at, then the Marathi for to grumble, then say \emph{sameṭ} again
\end{itemize}

\begin{tcolorbox}[breakable,colback=teal!4,colframe=teal!35!black,title={Before you move on}]
What does \mr{समेट} mean? (\textquotedblleft{}A reconciliation\textquotedblright{}.) Say it once more.
\end{tcolorbox}
\section[\texorpdfstring{\mr{गैरसमज} (gairsamaj)}{गैरसमज (gairsamaj)}]{\mr{गैरसमज} (gairsamaj) — a misunderstanding}

\label{lesson:MR-C454-gairsamaj}



\begin{quote}
Before the new one: say the Marathi for to grumble, then the Marathi for a reconciliation.
\end{quote}

\subsection*{You'll want to know: \mr{गैरसमज}}
\textbf{\mr{गैरसमज}} — \emph{gairsamaj} — \textquotedblleft{}a misunderstanding\textquotedblright{}.

\begin{quote}
\mr{एका} \mr{लहान} \mr{गैरसमजामुळे} \mr{दोन} \mr{मित्र} \mr{भांडले}.
\end{quote}

\begin{quote}
\emph{ekā lahān gairsamjāmuḷe don mitra bhāṇḍle.}
\end{quote}

\begin{quote}
Two friends quarrelled over a small misunderstanding.
\end{quote}

Masculine.

\mr{गैरसमज} \mr{होणे}, to misunderstand: \mr{माझा} \mr{गैरसमज} \mr{झाला}, I got it wrong.

\begin{cousinweb}
\textbf{\mr{गैरसमज}} (\emph{gairsamaj}) is a Marathi compound built from \emph{gair-}, \textquotedblleft{}not, other\textquotedblright{}, a loan from Persian \emph{ghair} usually traced to Arabic, and \mr{समज}, \textquotedblleft{}understanding\textquotedblright{}, the noun of \mr{समजणे}.
\end{cousinweb}

\begin{grammarlens}[title={What you've built}]
One more word for describing what you lived and felt.
\end{grammarlens}

\subsection*{Your turn}
\begin{itemize}
\raggedright
  \item \emph{Say it:} \emph{gairsamaj}
  \item \emph{Say it:} \emph{gairsamaj}, once more
  \item \emph{Say it:} \emph{ekā lahān gairsamjāmuḷe don mitra bhāṇḍle.} — Two friends quarrelled over a small misunderstanding.
  \item \emph{Recall:} say the Marathi for to grumble, then the Marathi for a reconciliation, then say \emph{gairsamaj} again
\end{itemize}

\begin{tcolorbox}[breakable,colback=teal!4,colframe=teal!35!black,title={Before you move on}]
What does \mr{गैरसमज} mean? (\textquotedblleft{}A misunderstanding\textquotedblright{}.) Say it once more.
\end{tcolorbox}
\section[\texorpdfstring{\mr{रुसणे} (rusṇe)}{रुसणे (rusṇe)}]{\mr{रुसणे} (rusṇe) — to sulk}

\label{lesson:MR-C454-rusne}



\begin{quote}
Before the new one: say the Marathi for a reconciliation, then the Marathi for a misunderstanding.
\end{quote}

\subsection*{You'll want to know: \mr{रुसणे}}
\textbf{\mr{रुसणे}} — \emph{rusṇe} — \textquotedblleft{}to sulk\textquotedblright{}.

\begin{quote}
\mr{खेळणं} \mr{मिळालं} \mr{नाही} \mr{म्हणून} \mr{बाळ} \mr{रुसलं}.
\end{quote}

\begin{quote}
\emph{kheḷṇaṁ miḷālaṁ nāhī mhaṇūn bāḷ ruslaṁ.}
\end{quote}

\begin{quote}
The baby sulked because it did not get the toy.
\end{quote}

Past (he …): \mr{तो} \mr{रुसला}, \emph{to ruslā}.

The noun \mr{रुसवा}, a sulk; \mr{रुसवा} \mr{काढणे}, to coax someone out of a sulk.

\begin{cousinweb}
\textbf{\mr{रुसणे}} (\emph{rusṇe}) is a tadbhava, inherited from Sanskrit \emph{ruṣyati}, \textquotedblleft{}is angry\textquotedblright{}; CDIAL lists Marathi \emph{rusṇẽ}, \textquotedblleft{}to be offended\textquotedblright{} (CDIAL 10794).
\end{cousinweb}

\begin{grammarlens}[title={What you've built}]
One more word for describing what you lived and felt.
\end{grammarlens}

\subsection*{Your turn}
\begin{itemize}
\raggedright
  \item \emph{Say it:} \emph{rusṇe}
  \item \emph{Say it:} \emph{rusṇe}, once more
  \item \emph{Say it:} \emph{kheḷṇaṁ miḷālaṁ nāhī mhaṇūn bāḷ ruslaṁ.} — The baby sulked because it did not get the toy.
  \item \emph{Recall:} say the Marathi for a reconciliation, then the Marathi for a misunderstanding, then say \emph{rusṇe} again
\end{itemize}

\begin{tcolorbox}[breakable,colback=teal!4,colframe=teal!35!black,title={Before you move on}]
What does \mr{रुसणे} mean? (\textquotedblleft{}To sulk\textquotedblright{}.) Say it once more.
\end{tcolorbox}
\section[\texorpdfstring{\mr{मनवणे} (manavṇe)}{मनवणे (manavṇe)}]{\mr{मनवणे} (manavṇe) — to coax round; to win over}

\label{lesson:MR-C454-manavne}



\begin{quote}
Before the new one: say the Marathi for a misunderstanding, then the Marathi for to sulk.
\end{quote}

\subsection*{You'll want to know: \mr{मनवणे}}
\textbf{\mr{मनवणे}} — \emph{manavṇe} — \textquotedblleft{}to coax round; to win over\textquotedblright{}.

\begin{quote}
\mr{रुसलेल्या} \mr{बहिणीला} \mr{त्याने} \mr{गोड} \mr{बोलून} \mr{मनवलं}.
\end{quote}

\begin{quote}
\emph{ruslelyā bahiṇīlā tyāne goḍ bolūn manavlaṁ.}
\end{quote}

\begin{quote}
He coaxed his sulking sister round with sweet words.
\end{quote}

Past, with -\mr{ने} (he …): \mr{त्याने} \mr{तिला} \mr{मनवलं}, \emph{tyāne tilā manavlaṁ}.

The person coaxed takes -\mr{ला}: \mr{त्याने} \mr{बहिणीला} \mr{मनवलं}.

\begin{cousinweb}
\textbf{\mr{मनवणे}} (\emph{manavṇe}) is inherited: CDIAL lists Marathi \emph{mānaviṇẽ}, \textquotedblleft{}to persuade\textquotedblright{}, under Sanskrit \emph{manyate}, \textquotedblleft{}thinks, honours\textquotedblright{} (CDIAL 9857); \mr{मनवणे} is the everyday form.
\end{cousinweb}

\begin{grammarlens}[title={What you've built}]
One more word for describing what you lived and felt.
\end{grammarlens}

\subsection*{Your turn}
\begin{itemize}
\raggedright
  \item \emph{Say it:} \emph{manavṇe}
  \item \emph{Say it:} \emph{manavṇe}, once more
  \item \emph{Say it:} \emph{ruslelyā bahiṇīlā tyāne goḍ bolūn manavlaṁ.} — He coaxed his sulking sister round with sweet words.
  \item \emph{Recall:} say the Marathi for a misunderstanding, then the Marathi for to sulk, then say \emph{manavṇe} again
\end{itemize}

\begin{tcolorbox}[breakable,colback=teal!4,colframe=teal!35!black,title={Before you move on}]
What does \mr{मनवणे} mean? (\textquotedblleft{}To coax round; to win over\textquotedblright{}.) Say it once more.
\end{tcolorbox}
\section[\texorpdfstring{\mr{हिरमुसणे} (hirmusṇe)}{हिरमुसणे (hirmusṇe)}]{\mr{हिरमुसणे} (hirmusṇe) — to be crestfallen}

\label{lesson:MR-C454-hirmusne}



\begin{quote}
Before the new one: say the Marathi for to sulk, then the Marathi for to coax round.
\end{quote}

\subsection*{You'll want to know: \mr{हिरमुसणे}}
\textbf{\mr{हिरमुसणे}} — \emph{hirmusṇe} — \textquotedblleft{}to be crestfallen\textquotedblright{}.

\begin{quote}
\mr{परीक्षेत} \mr{कमी} \mr{गुण} \mr{मिळाले} \mr{म्हणून} \mr{ती} \mr{हिरमुसली}.
\end{quote}

\begin{quote}
\emph{parīkṣet kamī guṇ miḷāle mhaṇūn tī hirmuslī.}
\end{quote}

\begin{quote}
She was crestfallen because she got low marks in the exam.
\end{quote}

Past (he …): \mr{तो} \mr{हिरमुसला}, \emph{to hirmuslā}.

Of a disappointed child or player: \mr{हिरमुसलेला} \mr{चेहरा}, a crestfallen face.

\begin{cousinweb}
\textbf{\mr{हिरमुसणे}} (\emph{hirmusṇe}) is a Marathi verb whose origin is not settled.
\end{cousinweb}

\begin{grammarlens}[title={What you've built}]
One more word for describing what you lived and felt.
\end{grammarlens}

\subsection*{Your turn}
\begin{itemize}
\raggedright
  \item \emph{Say it:} \emph{hirmusṇe}
  \item \emph{Say it:} \emph{hirmusṇe}, once more
  \item \emph{Say it:} \emph{parīkṣet kamī guṇ miḷāle mhaṇūn tī hirmuslī.} — She was crestfallen because she got low marks in the exam.
  \item \emph{Recall:} say the Marathi for to sulk, then the Marathi for to coax round, then say \emph{hirmusṇe} again
\end{itemize}

\begin{tcolorbox}[breakable,colback=teal!4,colframe=teal!35!black,title={Before you move on}]
What does \mr{हिरमुसणे} mean? (\textquotedblleft{}To be crestfallen\textquotedblright{}.) Say it once more.
\end{tcolorbox}
